\documentclass[11pt]{article}
\usepackage[letterpaper,margin=.68in,headheight=16pt,headsep=.18in,footskip=.28in]{geometry}
\usepackage{amsmath,array,enumitem,fancyhdr,hyperref}
\pagestyle{fancy}\fancyhf{}\renewcommand{\headrulewidth}{.3pt}
\fancyhead[L]{\small Bellingham Math Circle / Week 12 / Return visits}
\fancyhead[R]{\small Adult guide}
\fancyfoot[L]{\small RV-12-FAC-v1 / Draft and unpiloted}\fancyfoot[R]{\small\thepage}
\setlength{\emergencystretch}{2em}
\setlength{\parindent}{0pt}\setlength{\parskip}{7pt}
\setlist[itemize]{leftmargin=1.3em,itemsep=3pt,topsep=4pt}
\hypersetup{hidelinks,pdftitle={Week 12 return-visit facilitator guide}}
\begin{document}
{\Large\bfseries Mathematical overview}\par\medskip
\textbf{Facts and limits.} Noncrossing perfect pairings of six fixed circle positions may be restricted to unlike labels. For cyclic words RBRBRB, RRRBBB and RRBRBB the counts are respectively 5, 1 and 2. A cyclic R/B word admits an unlike noncrossing pairing exactly when its two label counts agree: remove an adjacent unlike pair, then continue.

A path with \(n\) up and \(n\) down steps that stays on or above its starting line is a Dyck path. Restricting its height to at most two leaves \(2^{n-1}\) paths for \(n\geq1\): each return-to-ground block is \(U(UD)^kD\), and block sizes are a composition of \(n\). For four pairs, height ceilings 1, 2, 3 and 4 allow 1, 8, 13 and 14 paths.

A legal adjacent swap DU to UD raises one valley by two height units. Repeated raising reaches the mountain \(U^nD^n\). Half the sum of intermediate-height differences from that mountain is the exact minimum number of swaps needed; every swap changes this score by one. Equal counts of steps alone do not permit a swap that goes below ground.

\textbf{Children's destination.} Distinguish existence from counting under a restriction, organize paths around their ground returns, and pair an actual shortest route with a lower bound.

\textbf{Entry map.} Problem 1 (student pp. 1--2): K--1 and up, oral rules, labels and pairing. Problems 2--3 (student p. 3): grades 2--3 and up, reading U/D or arrow cards and keeping height nonnegative. Problem 4 (student p. 4): grades 4--5, local changes and a picture-based optimality argument; no formula is required.

\textbf{New versus base.} Base pages enumerate ordinary pairings and paths and translate trees/codes. These visits add label compatibility, a height ceiling, and a move-distance question.
\newpage {\Large\bfseries Prepare a return visit}\par\medskip
\textbf{Status and use.} Three investigations extend one familiar theme. These companions are separate from the base packets and may support several visits. All pages are new and unpiloted. Printed labels describe prerequisites, not age restrictions. Give one page at a time; completing every page is not the goal.

\textbf{Materials and preparation.} For each pair, three R and three B lettered counters, pencil and eraser; short strings are optional because pencil chords make crossings visible. Provide five U and five D step cards, ruler and spare grid paper. Reserve six sets for eleven children (upper third child rotates as checker). A pairing circle should be 3–4 inches across; path grids use equal horizontal and vertical scale. Neither counting task requires exact tile fit.

\textbf{Staffing.} Current context: eleven children, fixed KK11 / 3333 / 445 tables, one adult anchored to each table. Use pairs; the upper third child rotates as reader or checker. Routine adult reading, arrow drawing or recording can leave the mathematical decisions with children. If adults must choose every move, use the earlier entry rather than run the investigation for them.

\textbf{Short common launch.} After free handling, put four labeled dots around a circle and join one unlike pair with string or pencil. Let another child check whether the remaining unlike pair crosses it. On a separate short walk, demonstrate one U and one D card from ground level and show the matching intermediate point. The pairing and step representations are familiar if the base packet was used; otherwise give its ordinary board demonstration before the page.

\textbf{Suggested first route.} Begin with unlike-color pairing attempts. Offer the ceiling-two path page only once children can act out a walk. Keep local swaps for a later visit if choosing and preserving a legal path still takes adult control.

\textbf{Flexible timing.} Allow 3--5 minutes of material play and 2--4 minutes for the common demonstration, then 15--25 minutes for one investigation and 5 minutes to share. A longer second visit can spend 30--45 minutes on the same investigation, including unfinished examples or its explanation. These are planning estimates, not observed timings. Do not require all three within one hour.

\textbf{Questions and hints.} Start with ``What can you try?'' and ``What would count as success?'' Check the actual rule before giving a strategy. Optional questions and hints are keyed below; use them only after a real attempt. A counter argument or drawing may be a complete explanation.

\textbf{Source and pilot record.} Mathematical examples and arguments here are original local follow-ups to the current base theme. Consulted teaching source: \emph{Math Circle by the Bay}, Preface pp. vii--x (local PDF pp. 8--11), on interaction, manipulatives, hints, explanations and themes spanning multiple years. Our exact launch, entry route and timing are design inferences. Year-1 \emph{Handouts 6}, Problems 6.2--6.6, provided a local example of connecting representations. Record actual problem, starting route, examples tried, what needed adult rescue, and what children want to resume. Distinguish observations from hypotheses. Counter fit, materials stock, procedures and classroom response have not been physically tested.
\newpage
\textbf{Problem 1: unlike partners.} Using fixed positions 1–6, the complete unlike-color pairings are:

RBRBRB: \(12|34|56,\ 12|36|45,\ 14|23|56,\ 16|23|45,\ 16|25|34\).

RRRBBB: \(16|25|34\).

RRBRBB: \(16|23|45,\ 16|25|34\).

Here \(16\) means join dots 1 and 6; bars separate pairs. These are adult compact records, not a new notation required of children. Draw the actual chords when checking their collection.

\textbf{Why a pairing always exists when counts agree.} Unless no dots remain, a balanced cyclic word has both labels and therefore two adjacent unlike dots. Join that pair near the boundary, remove it, and pair the remaining cyclic word. Removing one R and one B preserves equality; the new chord cannot cross any later chord. Reinsert the removed adjacent pairs in reverse order. Conversely, every unlike pair consumes one of each label, so unequal counts make a complete pairing impossible.

\textbf{Completeness of the small lists.} A chord from dot 1 divides the remaining dots into two sides, each containing a balanced number of labels and an even number of dots. Check all possible partners of 1, then pair each side independently. This is an optional hint after genuine attempts.

\textbf{Extension.} Invent a balanced label arrangement with as few pairings as possible and one with many. Treat maximum counts on larger circles as a new question, not a theorem inferred from the six-dot trials.
\newpage
\textbf{Problems 2--3: a low ceiling.} With four U and four D cards, the eight paths staying at height at most two are
\[
\begin{array}{ll}
UDUDUDUD&UUDDUDUD\\
UDUUDDUD&UDUDUUDD\\
UUDDUUDD&UUDUDDUD\\
UDUUDUDD&UUDUDUDD.
\end{array}
\]
Problem 3 has \textbf{sixteen} five-pair paths under the same ceiling: the four possible gaps between pair tokens may each be a block break or not. This continues the same bounded-height investigation rather than adding a fourth investigation.

A ceiling of one allows only UDUDUDUD. A ceiling of three permits every ordinary four-pair path except UUUUDDDD, giving 13; a ceiling of four permits all 14.

\textbf{Why eight.} Between consecutive ground returns, a ceiling-two path must begin U, remain at height one or two by repeated UD, then end D. A block of size \(r\) is therefore \(U(UD)^{r-1}D\). Divide four pairs into an ordered list of positive block sizes: \(1111,211,121,112,22,31,13,4\). Each gives exactly one path. For \(n\) pairs there are \(n-1\) gaps between pair tokens; choose independently whether each gap separates blocks, giving \(2^{n-1}\). This is a deeper adult continuation, not a printed procedure.

\textbf{If needed.} Use a pawn to follow one path before asking for a collection. Put a ruler at the ceiling; stop attempts immediately when they cross it. Children may organize their own list before an adult suggests ground returns.

\textbf{Return question.} What happens at a ceiling of three with more pairs? Concrete enumeration is welcome; the simple ceiling-two formula should not be transferred.
\newpage
\textbf{Problem 4: raise the path.} A swap DU to UD changes only the height at their intermediate vertex, raising it by two. It always preserves legality. The reverse swap is legal only if its new intermediate height is nonnegative.

The four-pair alternating path reaches UUUUDDDD in six moves. One shortest route is
\[
\begin{array}{l}
UDUDUDUD\to UUDDUDUD\to UUDUDDUD\to UUUDDDUD\\
\to UUUDDUDD\to UUUDUDDD\to UUUUDDDD.
\end{array}
\]
Each arrow swaps adjacent unlike letters and leaves a legal path. The other printed starts UUDDUUDD and UUUDUDDD need respectively four moves and one. A shortest route from UUDDUUDD is UUDDUUDD $\to$ UUDUDUDD $\to$ UUUDDUDD $\to$ UUUDUDDD $\to$ UUUUDDDD. UUUDUDDD becomes the target by swapping its middle DU.

\textbf{Why no shorter route.} Above the alternating path and below the mountain, there are six elementary height-increment diamonds. Equivalently, intermediate-height differences sum to twelve, and one swap can change that sum by only two. Therefore at least six swaps are necessary. Draw the two paths on the same grid and let children account for the six raised pieces; avoid demanding a concurrent area tally during the route.

\textbf{General reason.} Any nonmountain Dyck path has a valley: some D occurs before a later U, so an adjacent DU exists. Raise it. Half the total height gap to the mountain falls by one and stays nonnegative. Thus raising terminates precisely at the mountain, with the lower-bound number of moves. This proves both connectivity and optimality to that target.

\textbf{If needed.} Supply a shorter legal path and demonstrate a single swap on separate cards. Keep the target question open. A score alone does not give the distance between two arbitrary paths; that extension requires a new argument.

\textbf{Check limit.} Exhaustive enumeration verifies all four-pair paths, ceilings and legal swaps. The finite test is separate from the general valley and height-gap proof.
\end{document}
